\documentclass[11pt]{beamer}
\usepackage[utf8]{inputenc}
\usepackage[T1,T2A]{fontenc}
\usepackage[russian]{babel}
\usepackage{color}
\usepackage{listings}
\usepackage{calc}
\usepackage{graphicx}
\usepackage{epstopdf}
\usepackage{hyperref}
\hypersetup{unicode,colorlinks}
\usepackage{csquotes}
\usepackage{upquote}
\usepackage{cprotect}
\usetheme[progressbar=head,numbering=fraction,block=fill]{metropolis}
\usepackage{dejavu}
\usepackage{etoolbox}
\usepackage{bm}
\usepackage{minted}
%\usepackage[defaultsans]{droidsans}
%\usepackage[defaultmono]{droidmono}
%\makeatletter
%\input droid/t1fdm.fd
%\makeatother

\makeatletter
\DeclareRobustCommand*{\bfseries}{%
    \not@math@alphabet\bfseries\mathbf
    \fontseries\bfdefault\selectfont
    \boldmath
}
\makeatother

\makeatletter
\chardef\straightquote@code=\catcode`'
\chardef\backquote@code=\catcode``
\catcode`'=\active \catcode``=\active
\patchcmd{\@noligs}
{\textasciigrave}
{\fixedtextasciigrave}
{}{}
\newcommand{\fixedtextasciigrave}{%
    \makebox[.5em]{\fontencoding{TS1}\fontfamily{fvs}\selectfont\textasciigrave}% Vera Sans
}
\catcode\lq\'=\straightquote@code
\catcode\lq\`=\backquote@code
\makeatletter

\lstset{
    language=Haskell,
    inputencoding=utf8,
    extendedchars=true,
    breaklines=true,
    escapeinside=!!,
    tabsize=4,
    breakatwhitespace=true,
    keepspaces=true,
    basicstyle=\ttfamily,
    commentstyle=\color{gray}\itshape,
    keywordstyle=
    % keywordstyle=\color{blue}, % TODO https://tex.stackexchange.com/questions/415777/avoid-highlighting-keywords-following-certain-words-in-listings
    % addkeywords={Type}
    % deletekeywords={Monad, String, return, array, show},
    % otherkeywords={::, ->}
}
\lstset{
    literate={а}{{\selectfont\char224}}1
    {б}{{\selectfont\char225}}1
    {в}{{\selectfont\char226}}1
    {г}{{\selectfont\char227}}1
    {д}{{\selectfont\char228}}1
    {е}{{\selectfont\char229}}1
    {ё}{{\"e}}1
    {ж}{{\selectfont\char230}}1
    {з}{{\selectfont\char231}}1
    {и}{{\selectfont\char232}}1
    {й}{{\selectfont\char233}}1
    {к}{{\selectfont\char234}}1
    {л}{{\selectfont\char235}}1
    {м}{{\selectfont\char236}}1
    {н}{{\selectfont\char237}}1
    {о}{{\selectfont\char238}}1
    {п}{{\selectfont\char239}}1
    {р}{{\selectfont\char240}}1
    {с}{{\selectfont\char241}}1
    {т}{{\selectfont\char242}}1
    {у}{{\selectfont\char243}}1
    {ф}{{\selectfont\char244}}1
    {х}{{\selectfont\char245}}1
    {ц}{{\selectfont\char246}}1
    {ч}{{\selectfont\char247}}1
    {ш}{{\selectfont\char248}}1
    {щ}{{\selectfont\char249}}1
    {ъ}{{\selectfont\char250}}1
    {ы}{{\selectfont\char251}}1
    {ь}{{\selectfont\char252}}1
    {э}{{\selectfont\char253}}1
    {ю}{{\selectfont\char254}}1
    {я}{{\selectfont\char255}}1
    {А}{{\selectfont\char192}}1
    {Б}{{\selectfont\char193}}1
    {В}{{\selectfont\char194}}1
    {Г}{{\selectfont\char195}}1
    {Д}{{\selectfont\char196}}1
    {Е}{{\selectfont\char197}}1
    {Ё}{{\"E}}1
    {Ж}{{\selectfont\char198}}1
    {З}{{\selectfont\char199}}1
    {И}{{\selectfont\char200}}1
    {Й}{{\selectfont\char201}}1
    {К}{{\selectfont\char202}}1
    {Л}{{\selectfont\char203}}1
    {М}{{\selectfont\char204}}1
    {Н}{{\selectfont\char205}}1
    {О}{{\selectfont\char206}}1
    {П}{{\selectfont\char207}}1
    {Р}{{\selectfont\char208}}1
    {С}{{\selectfont\char209}}1
    {Т}{{\selectfont\char210}}1
    {У}{{\selectfont\char211}}1
    {Ф}{{\selectfont\char212}}1
    {Х}{{\selectfont\char213}}1
    {Ц}{{\selectfont\char214}}1
    {Ч}{{\selectfont\char215}}1
    {Ш}{{\selectfont\char216}}1
    {Щ}{{\selectfont\char217}}1
    {Ъ}{{\selectfont\char218}}1
    {Ы}{{\selectfont\char219}}1
    {Ь}{{\selectfont\char220}}1
    {Э}{{\selectfont\char221}}1
    {Ю}{{\selectfont\char222}}1
    {Я}{{\selectfont\char223}}1
    {`}{{{\fixedtextasciigrave}}}1
}

% \usemintedstyle{solarized-light}
\newminted[haskell]{haskell}{
    escapeinside=``,
    beameroverlays=true,
    texcomments=true,
    autogobble=true,
    fontsize=\small
}
\newmintinline[haskinline]{haskell}{
    escapeinside=``,
    beameroverlays=true,
    texcomments=true
}

\author{Алексей Романов}
\subtitle{Функциональное программирование на Haskell}
%\logo{}
\institute{МИЭТ}
\subject{Функциональное программирование на Haskell}
%\setbeamercovered{transparent}
%\setbeamertemplate{navigation symbols}{}

\setminted[lean4]{extrakeywords={rfl,rw,change,induction}}

\title{Лекция 14: рассуждения с помощью равенств и структурная индукция}

\begin{document}
\begin{frame}[plain]
  \maketitle
\end{frame}

\begin{frame}[fragile]
  \frametitle{От проверки свойства к доказательству}
  Пусть \haskinline|xs| и \haskinline|ys| "--- списки целых чисел. Как проверить, что порядок двух операций не влияет на результат?
  \begin{haskellsmall}
    prop_mapAppend = property $ do
      xs <- forAll genList
      ys <- forAll genList
      map (+1) (xs ++ ys) === map (+1) xs ++ map (+1) ys
  \end{haskellsmall}
  Здесь \haskinline|genList :: Gen [Integer]|. Допустим, все 100 тестов прошли. Что мы узнали? \pause

  Равенство выполнено на этих 100 парах. Можно ли доказать его сразу для всех конечных списков и заменить \haskinline|(+1)| произвольной функцией?
\end{frame}

\begin{frame}[fragile]
  \frametitle{Какие значения мы рассматриваем?}
  \begin{itemize}
    \item В этой лекции списки и деревья конечны и полностью определены. Функции на рассматриваемых входах возвращают полностью определённые значения.
    \item Зачем это оговаривать? Список \haskinline|repeat 1| тоже имеет тип \haskinline|[Integer]|, но до его конца мы никогда не дойдём. \pause
    \item В доказательстве $=$ означает равенство значений. Его можно доказывать и для типов, у которых нет экземпляра \haskinline|Eq|.
  \end{itemize}
\end{frame}

\begin{frame}[fragile]
  \frametitle{Определение как равенство}
  \begin{haskell}
    []       ++ ys = ys
    (x : xs) ++ ys = x : (xs ++ ys)

    map f []       = []
    map f (x : xs) = f x : map f xs
  \end{haskell}
  Если образец подходит, одну сторону уравнения можно заменить другой. Например, что получится из \haskinline|map f (x : [])|? \pause
  \begin{haskell}
      map f (x : [])
    !$=$! f x : map f []
    !$=$! f x : []
  \end{haskell}
  Во втором переходе заменили выражение внутри \haskinline|(:)|.
\end{frame}

\begin{frame}[fragile]
  \frametitle{Где вычисление останавливается?}
  Хотим доказать для любых \haskinline|f|, \haskinline|xs| и \haskinline|ys|:
  \begin{haskell}
    map f (xs ++ ys) !$=$! map f xs ++ map f ys
  \end{haskell}
  Какое определение можно раскрыть слева? \pause
  \begin{itemize}
    \item Чтобы раскрыть \haskinline|(++)|, нужно знать первый конструктор \haskinline|xs|. Вид списка \haskinline|ys| здесь не важен.
    \item Зафиксируем произвольные \haskinline|f| и \haskinline|ys|. Докажем утверждение $P(\mathit{xs})$ индукцией по \haskinline|xs|.
  \end{itemize}
\end{frame}

\begin{frame}[fragile]
  \frametitle{Структурная индукция по списку}
  \begin{itemize}
    \item \emph{База.} Докажите $P([])$.
    \item \emph{Шаг.} Возьмите произвольные голову \haskinline|x| и хвост \haskinline|xs|. Предположив $P(\mathit{xs})$, докажите $P(x : \mathit{xs})$.
    \item Почему этого достаточно для любого конечного списка? \pause
    \item Например, для \haskinline|[a,b]| база даёт $P([])$. Шаг с головой \haskinline|b| даёт $P([b])$, ещё один шаг с головой \haskinline|a| "--- $P([a,b])$.
  \end{itemize}
\end{frame}

\begin{frame}[fragile]
  \frametitle{База: первый список пуст}
  Подставим \haskinline|[]| вместо \haskinline|xs|. Как упростить обе стороны? \pause
  \begin{haskell}
      map f ([] ++ ys)
    !$=$! map f ys
  \end{haskell}
  \pause
  \begin{haskell}
      map f [] ++ map f ys
    !$=$! [] ++ map f ys
    !$=$! map f ys
  \end{haskell}
  Получили одно и то же выражение при любом \haskinline|ys|.
\end{frame}

\begin{frame}[fragile]
  \frametitle{Шаг: где применить предположение?}
  Для хвоста уже предполагаем:
  \begin{haskellsmall}
    map f (xs ++ ys) !$=$! map f xs ++ map f ys
  \end{haskellsmall}
  Преобразуем левую сторону для \haskinline|x : xs|:
  \begin{haskellsmall}
      map f ((x : xs) ++ ys)
    !$=$! map f (x : (xs ++ ys))
    !$=$! f x : map f (xs ++ ys)
  \end{haskellsmall}
  \pause
  \begin{haskellsmall}
    !$=$! f x : (map f xs ++ map f ys)
    !$=$! (f x : map f xs) ++ map f ys
    !$=$! map f (x : xs) ++ map f ys
  \end{haskellsmall}
  В первом из этих переходов применили предположение к хвосту. В двух последних прочитали определения справа налево.
\end{frame}

\begin{frame}[fragile]
  \frametitle{Lean: как читать доказательство?}
  В \href{https://lean-lang.org/}{Lean} \href{https://homepages.inf.ed.ac.uk/wadler/papers/propositions-as-types/propositions-as-types.pdf}{утверждение — это тип}, а его доказательство — значение этого типа.
  \begin{minted}[autogobble,fontsize=\footnotesize,escapeinside=!!]{lean4}
    example {A : Type} (xs : List A) : [] ++ xs = xs := by
      !\pause!rfl
  \end{minted}
  \begin{itemize}
    \item \mintinline{lean4}{example} задаёт утверждение для любого списка. После \mintinline{lean4}{by} идут шаги доказательства.
    \item \mintinline{lean4}|{A : Type}| "--- неявный аргумент: Lean выводит тип элемента из типа списка. \mintinline{lean4}|(xs : List A)| "--- явный аргумент: список передаём сами.
    \item По определению конкатенации левая часть сводится к \mintinline{lean4}{xs}. Команда \mintinline{lean4}{rfl} (reflexivity) закрывает цель, когда обе стороны совпадают по определению.
  \end{itemize}
\end{frame}

\begin{frame}[fragile]
  \frametitle{Тот же шаг индукции в Lean 4}
  Запишем цель и выберем индукцию по первому списку.
  \begin{minted}[autogobble,fontsize=\footnotesize,escapeinside=!!]{lean4}
    example {A B : Type} (f : A -> B) (xs ys : List A) :
        (xs ++ ys).map f = xs.map f ++ ys.map f := by
      induction xs with
      | nil => rfl
      | cons x xs ih =>
        change f x :: (xs ++ ys).map f =
               f x :: (xs.map f ++ ys.map f)
        !\pause!rw [ih]
  \end{minted}
  В базе \mintinline{lean4}{rfl} проверяет равенство по определению.
  В шаге \mintinline{lean4}{change} раскрывает определения, а \mintinline{lean4}{rw [ih]} (rewrite) переписывает левую часть по предположению для хвоста и закрывает цель.
\end{frame}

\begin{frame}[fragile]
  \frametitle{Liquid Haskell: уточнённые типы}
  \href{https://ucsd-progsys.github.io/liquidhaskell/}{Liquid Haskell} позволяет дополнить тип условием на его значения. Например, \texttt{\{v:Integer | v > 0\}} "--- тип положительных целых чисел.
  \begin{minted}[autogobble,fontsize=\footnotesize]{haskell}
    {-@ inc :: x:Integer -> {v:Integer | v > x} @-}
    inc :: Integer -> Integer
    inc x = x + 1
  \end{minted}
  Здесь условие на результат \haskinline|v| зависит от аргумента \haskinline|x|. Подойдёт ли реализация \haskinline|inc x = x|? \pause

  Она пройдёт проверку обычного типа. Liquid Haskell отклонит её: возвращённое значение должно быть строго больше аргумента.
\end{frame}

\begin{frame}[fragile]
  \frametitle{Какие определения раскрывает Liquid Haskell?}
  \begin{minted}[autogobble,fontsize=\footnotesize]{haskell}
    {-@ LIQUID "--reflection" @-}
    import Prelude hiding (map, (++))

    {-@ reflect map @-}
    map :: (a -> b) -> [a] -> [b]
    map _ []       = []
    map f (x : xs) = f x : map f xs

    {-@ infixr 5 ++ @-}
    {-@ reflect ++ @-}
    (++) :: [a] -> [a] -> [a]
    []       ++ ys = ys
    (x : xs) ++ ys = x : (xs ++ ys)
  \end{minted}
  \texttt{reflect} позволяет использовать уравнения этих функций в доказательстве.
\end{frame}

\begin{frame}[fragile]
  \frametitle{Liquid Haskell: база индукции без PLE}
  Доказываемое равенство записываем в уточнении типа результата.
  \href{https://github.com/ucsd-progsys/liquidhaskell/blob/master/liquid-prelude/src/Language/Haskell/Liquid/ProofCombinators.hs}{Документация в комментариях}.
  \begin{minted}[autogobble,fontsize=\footnotesize]{haskell}
    import Language.Haskell.Liquid.ProofCombinators

    {-@ mapAppend :: f:(a -> b) -> xs:[a] -> ys:[a]
          -> { map f (xs ++ ys) == map f xs ++ map f ys } @-}
    mapAppend :: (a -> b) -> [a] -> [a] -> Proof
    mapAppend f [] ys =
          map f ([] ++ ys)
      === map f ys
      === ([] ++ map f ys)
      === (map f [] ++ map f ys)
      *** QED
  \end{minted}
  Каждый переход \haskinline|(===)| требует доказать равенство соседних выражений. \haskinline|*** QED| завершает цепочку.
\end{frame}

\begin{frame}[fragile]
  \frametitle{Liquid Haskell: шаг индукции без PLE}
  Как получить равенство для хвоста \haskinline|xs|?
  \pause
  \begin{minted}[autogobble,fontsize=\footnotesize]{haskell}
    mapAppend f (x : xs) ys =
          map f ((x : xs) ++ ys)
      === map f (x : (xs ++ ys))
      === (f x : map f (xs ++ ys))
        ? mapAppend f xs ys
      === (f x : (map f xs ++ map f ys))
      === ((f x : map f xs) ++ map f ys)
      === (map f (x : xs) ++ map f ys)
      *** QED
  \end{minted}
  Подсказка \haskinline|? mapAppend f xs ys| позволяет заменить хвост под \haskinline|(:)|.
  \pause
  \par Те же шаги, что и в доказательстве на бумаге. Liquid Haskell автоматически проверяет каждый переход.
\end{frame}

\begin{frame}[fragile]
  \frametitle{То же доказательство с PLE}
  PLE (\emph{Proof by Logical Evaluation}) автоматически раскрывает определения \haskinline|map| и \haskinline|(++)|.
  \begin{minted}[autogobble,fontsize=\footnotesize]{haskell}
    {-@ LIQUID "--ple" @-}  -- включаем PLE для всего модуля

    mapAppend _ []       _  = ()
  \end{minted}
  Что остаётся от шага индукции?
  \begin{haskellsmall}
    mapAppend f (_ : xs) ys =!\pause! mapAppend f xs ys
  \end{haskellsmall}
  Рекурсивный вызов даёт предположение индукции, как \mintinline{lean4}{rw [ih]} в Lean.
  \pause\par Все остальные шаги доказательства PLE выполняет автоматически. При желании любой из них можно записать явно.
\end{frame}

\begin{frame}[fragile]
  \frametitle{Упражнение: ассоциативность}
  Докажите закон, знакомый по списковому моноиду:
  \begin{haskell}
    (xs ++ ys) ++ zs !$=$! xs ++ (ys ++ zs)
  \end{haskell}
  \begin{itemize}
    \item По какому списку провести индукцию? Какие два списка можно зафиксировать?
    \item Выпишите предположение для хвоста и отметьте ровно тот переход, в котором оно понадобилось.
  \end{itemize}
\end{frame}

\begin{frame}[fragile]
  \frametitle{Деревья: сколько предположений?}
  Вернёмся к дереву из лекции 6:
  \begin{haskell}
    data Tree a = Leaf a | Node (Tree a) (Tree a)

    mirror (Leaf x)   = Leaf x
    mirror (Node l r) = Node (mirror r) (mirror l)
  \end{haskell}
  Хотим доказать \haskinline|mirror (mirror t)| $=$ \haskinline|t|. Какие случаи придётся разобрать? \pause
  \begin{itemize}
    \item Для \haskinline|Leaf x| достаточно раскрыть определение дважды.
    \item Для \haskinline|Node l r| есть два предположения: одно для \haskinline|l|, другое для \haskinline|r|.
  \end{itemize}
\end{frame}

\begin{frame}[fragile]
  \frametitle{Два отражения узла}
  \begin{haskell}
      mirror (mirror (Node l r))
    !$=$! mirror (Node (mirror r) (mirror l))
    !$=$! Node (mirror (mirror l)) (mirror (mirror r))
  \end{haskell}
  Какое подвыражение можно заменить, зная свойство только для \haskinline|l|? \pause
  \begin{haskell}
    !$=$! Node l (mirror (mirror r))
  \end{haskell}
  \pause
  \begin{haskell}
    !$=$! Node l r
  \end{haskell}
  В каждом предположении фигурирует своё непосредственное поддерево.
\end{frame}

\begin{frame}[fragile]
  \frametitle{Когда предположение слишком слабое}
  Зададим обращение списка и сравним его с версией с аккумулятором:
  \begin{haskellsmall}
    reverse []       = []
    reverse (x : xs) = reverse xs ++ [x]

    revAcc []       acc = acc
    revAcc (x : xs) acc = revAcc xs (x : acc)
  \end{haskellsmall}
  Хотим доказать \haskinline|revAcc xs []| $=$ \haskinline|reverse xs|. Что произойдёт на первом шаге? \pause
  \begin{haskell}
      revAcc (x : xs) []
    !$=$! revAcc xs [x]
  \end{haskell}
  Предположение говорит только о вызове с \haskinline|[]|. Какое свойство охватит и новый аккумулятор?
\end{frame}

\begin{frame}[fragile]
  \frametitle{Усилим утверждение}
  Для каждого \haskinline|xs| докажем равенство \emph{при любом} \haskinline|acc|:
  \begin{haskell}
    revAcc xs acc !$=$! reverse xs ++ acc
  \end{haskell}
  База следует из определения. В шаге выберем для хвоста аккумулятор \haskinline|x : acc|: \pause
  \begin{haskellsmall}
      revAcc (x : xs) acc
    !$=$! revAcc xs (x : acc)
    !$=$! reverse xs ++ (x : acc)
    !$=$! reverse xs ++ ([x] ++ acc)
    !$=$! (reverse xs ++ [x]) ++ acc
    !$=$! reverse (x : xs) ++ acc
  \end{haskellsmall}
  Использовали ассоциативность и уравнение для \haskinline|reverse|. При \haskinline|acc = []| получаем исходное свойство по закону правой единицы.
\end{frame}

\begin{frame}[fragile]
  \frametitle{Для самостоятельной работы}
  Задания 1 и 2: попробуйте доказать вручную, затем в Lean и Liquid Haskell.
  \begin{enumerate}
    \item Докажите закон правой единицы \haskinline|xs ++ []| $=$ \haskinline|xs|. Почему для левой единицы индукция не понадобилась?
    \item Для \haskinline|Tree Integer| докажите: \haskinline|sumTree (mirror t)| $=$ \haskinline|sumTree t|. Какое свойство сложения понадобится в случае \haskinline|Node|?
    \item Всегда ли отражение сохраняет результат \haskinline|foldTree id op|? Найдите контрпример. Какое свойство \haskinline|op| гарантирует равенство результатов?
  \end{enumerate}
\end{frame}

\begin{frame}
  \frametitle{Дополнительное чтение}
  \small
  \begin{itemize}
    \setlength{\itemsep}{1pt}
    \item \href{https://nikivazou.github.io/floc26/Lecture_01_RefinementTypes.html}{Refinement Types} (Niki Vazou), \href{https://ruhaskell.org/posts/utils/2016/12/16/liquidhaskell-hello.html}{LiquidHaskell: знакомство} и \href{https://liquid.kosmikus.org/}{Liquid Haskell Tutorial}.
    \item \href{https://doi.org/10.1145/3242744.3242756}{Theorem Proving for All}, \S\,2.3: индукция по списку в Liquid Haskell.
    \item Graham Hutton, \href{https://people.cs.nott.ac.uk/pszgmh/pih.html}{Programming in Haskell} (платная книга), гл. 16 о доказательствах.
    \item Graham Hutton, \href{https://people.cs.nott.ac.uk/pszgmh/afp.html}{видеолекции 11--14}: обращение списка и обход дерева.
    \item \href{https://fstar-lang.org/tutorial/book/}{Proof-Oriented Programming in F*} (бесплатная книга). \href{https://fstar-lang.org/tutorial/book/part1/part1_lemmas.html}{Доказательство свойств рекурсивных функций по индукции}.
    \item \href{https://logicforprogrammers.com/}{Logic for Programmers} (платная книга), гл. 2--5: от равенств до Dafny.
  \end{itemize}
\end{frame}
\begin{frame}
  \frametitle{Дополнительное чтение: формальные методы}
  \small
  \begin{itemize}
    \item \href{https://www.hillelwayne.com/post/theorem-prover-showdown/}{The Great Theorem Prover Showdown}. Три задачи на проверку программ и их решения в разных системах.
    \item \href{https://www.hillelwayne.com/post/lpl/}{Let's Prove Leftpad}. Продолжение: как доказать корректность одной функции в разных системах.
    \item \href{https://www.cs.ru.nl/~freek/100/}{Formalizing 100 Theorems}. Иррациональность $\sqrt{2}$ и другие теоремы в разных системах.
    \item \href{https://blog.acolyer.org/2015/01/12/how-to-write-a-21st-century-proof/}{How to Write a 21st Century Proof}. Обзор статьи Лампорта о том, как разбить доказательство на проверяемые шаги.
  \end{itemize}
\end{frame}
\begin{frame}
  \frametitle{Дополнительное чтение: Lean}
  \small
  \begin{itemize}
    \setlength{\itemsep}{1pt}
    \item \href{https://lean-lang.org/functional_programming_in_lean/}{Functional Programming in Lean}. Введение для программистов.
    \item \href{https://lean-lang.org/theorem_proving_in_lean4/}{Theorem Proving in Lean 4}. От \href{https://lean-lang.org/theorem_proving_in_lean4/Introduction/}{введения} к тактикам и доказательствам по индукции.
    \item \href{https://leanprover-community.github.io/mathematics_in_lean/}{Mathematics in Lean}. Математические доказательства с библиотекой mathlib.
    \item \href{https://adam.math.hhu.de/}{Lean Game Server}. Интерактивные задачи на доказательство.
    \item \href{https://leanprover-community.github.io/logic_and_proof/}{Logic and Proof}, главы 17--18. Индукция на бумаге и в Lean 4.
    \item Heather Macbeth, \href{https://hrmacbeth.github.io/math2001/}{The Mechanics of Proof}, главы 1 и 6. Равенства и индукция в Lean 4.
    \item \href{https://leanprover-community.github.io/}{Lean community}: учебные материалы и сообщество Lean.
  \end{itemize}
\end{frame}
\end{document}